\documentclass[11pt,a4paper]{article}

\usepackage[margin=2.5cm, headheight=14pt]{geometry}
\usepackage{booktabs}
\usepackage{multirow}
\usepackage{array}
\usepackage{xcolor}
\usepackage{hyperref}
\usepackage{amsmath}
\usepackage{microtype}
\usepackage{parskip}
\usepackage{enumitem}
\usepackage{caption}
\usepackage{subcaption}
\usepackage{fancyhdr}
\usepackage{titlesec}
\usepackage{tabularx}
\usepackage{colortbl}

% Colour accents
\definecolor{ownerblue}{RGB}{30,80,160}
\definecolor{lightgray}{RGB}{245,245,245}
\definecolor{bertcolor}{RGB}{220,240,255}
\definecolor{llmcolor}{RGB}{220,255,230}
\definecolor{wincolor}{RGB}{180,220,180}

\hypersetup{
  colorlinks=true,
  linkcolor=ownerblue,
  urlcolor=ownerblue,
  citecolor=ownerblue,
}

\pagestyle{fancy}
\fancyhf{}
\rhead{\textcolor{gray}{\small OWNER -- Cluster Naming Evaluation}}
\lhead{\textcolor{gray}{\small BERTScore Report}}
\rfoot{\thepage}

\titleformat{\section}{\large\bfseries\color{ownerblue}}{{\thesection}}{1em}{}
\titleformat{\subsection}{\normalsize\bfseries}{{\thesubsection}}{1em}{}

\begin{document}

% ─── Title ───────────────────────────────────────────────────────────────────
\begin{center}
  {\LARGE\bfseries\color{ownerblue}
    Cluster Naming Evaluation Report\\[0.4em]
    \large OWNER: Toward Unsupervised Open-World Named Entity Recognition}\\[1em]
  {\large Comparison of BERT-based and LLM-based (Llama-3.1-8B) Naming\\
  Evaluated with BERTScore across Five CrossNER Domains}\\[0.8em]
  {\small \today}
\end{center}
\vspace{0.5em}
\hrule
\vspace{1.5em}

% ─── Abstract ────────────────────────────────────────────────────────────────
\begin{abstract}
OWNER~\cite{genest2025owner} organises extracted entities into clusters
whose integer IDs carry no semantic meaning.
This report evaluates two automatic cluster-naming strategies proposed in the
paper: (1)~\textbf{BERT-based naming}, which predicts the most frequent
\texttt{[MASK]} token across all entities in a cluster, and
(2)~\textbf{LLM-based naming} (Llama-3.1-8B via Ollama), which infers a type
label from a sample of $n{=}16$ entities using the Figure~3 prompt.
Named clusters are scored against ground-truth entity type labels using
\textbf{BERTScore} (RoBERTa-large backbone) across five CrossNER domains:
AI, Literature, Music, Politics, and Science.
BERT naming achieves an average F$_1$ of \textbf{0.966} while LLM naming
achieves \textbf{0.944}, with BERT performing consistently better except
on the Music domain where both methods are within 0.001 F$_1$ of each other.
\end{abstract}

% ─── 1. Introduction ─────────────────────────────────────────────────────────
\section{Introduction}

Named Entity Recognition (NER) traditionally requires a fixed, predefined
set of entity types. OWNER removes this constraint by clustering entity
embeddings into semantically coherent groups, automatically discovering the
number and nature of entity types without annotation. However, the resulting
clusters are identified only by integer IDs. Assigning meaningful names to
these clusters is necessary for interpretability, knowledge graph construction,
and human evaluation.

The paper proposes two complementary naming strategies (Section~III-B-3-C):

\begin{itemize}[leftmargin=1.5em]
  \item \textbf{BERT naming}: leverage the same BERT encoder used for entity
    typing. For each entity, the prompt
    $\mathcal{P}(e,X)=\texttt{``\{sentence\} \{entity\} is a [MASK].''}$
    is passed through the MLM head; the top-1 predicted token is collected.
    The most frequent predicted token across all entities in a cluster
    becomes its name. This approach is cost-free as the encoder is already
    available.
  \item \textbf{LLM naming}: a random sample of $n=16$ entity strings is
    presented to Llama-3.1-8B-Instruct (via Ollama) using the prompt in
    Figure~3 of the paper (temperature~$= 0$). The model responds with a
    single entity type name.
\end{itemize}

\noindent
Both approaches are evaluated against ground-truth labels using
BERTScore~\cite{zhang2020bertscore}, which measures token-level semantic
similarity via contextualised RoBERTa-large embeddings.

% ─── 2. Experimental Setup ───────────────────────────────────────────────────
\section{Experimental Setup}

\subsection{Datasets}

We evaluate on five CrossNER target domains used in the OWNER paper.
Table~\ref{tab:datasets} summarises the number of entity types and aligned
clusters in each domain.

\begin{table}[h]
\centering
\caption{CrossNER target domains and cluster statistics.}
\label{tab:datasets}
\setlength{\tabcolsep}{8pt}
\begin{tabular}{l r r}
\toprule
\textbf{Domain} & \textbf{True entity types} $k$ & \textbf{Aligned clusters} \\
\midrule
AI         & 14 & 16 \\
Literature & 12 & 14 \\
Music      & 13 & 20 \\
Politics   & 9  & 22 \\
Science    & 17 & 20 \\
\bottomrule
\end{tabular}
\end{table}

\subsection{Ground-Truth Alignment}

The ground-truth label for each cluster is determined by finding the
true entity type that contributes the highest proportion of entities to
that cluster (column-wise argmax of the confusion matrix), stored as
\texttt{et\_test\_cluster\_naming\_gt.json}.
Only clusters present in all three files (GT, BERT, LLM) are included in
the evaluation.

\subsection{BERTScore}

BERTScore computes precision, recall, and F$_1$ by greedily matching
tokens of the candidate string to tokens of the reference string using
cosine similarity of RoBERTa-large contextual embeddings.
We report macro-averaged P, R, F$_1$ across all aligned clusters within
each domain, and then macro-average over domains.

% ─── 3. Results ──────────────────────────────────────────────────────────────
\section{Results}

\subsection{BERTScore per Domain}

Table~\ref{tab:bertscore} presents the full BERTScore results.
Bold entries indicate the better F$_1$ for each domain.

\begin{table}[h]
\centering
\caption{BERTScore (Precision / Recall / F$_1$) of BERT-based and
LLM-based naming against ground-truth labels.
$k$ = number of aligned clusters.
Best F$_1$ per domain in \textbf{bold}.}
\label{tab:bertscore}
\setlength{\tabcolsep}{5pt}
\begin{tabular}{l r ccc r ccc}
\toprule
& \multicolumn{4}{c}{\cellcolor{bertcolor}\textbf{BERT Naming}}
& \multicolumn{4}{c}{\cellcolor{llmcolor}\textbf{LLM Naming (Llama-3.1-8B)}} \\
\cmidrule(lr){2-5}\cmidrule(lr){6-9}
\textbf{Domain} & $k$ & P & R & F$_1$ & $k$ & P & R & F$_1$ \\
\midrule
AI         & 16 & 0.992 & 0.986 & \textbf{0.989} & 16 & 0.915 & 0.945 & 0.929 \\
Literature & 14 & 0.974 & 0.975 & \textbf{0.974} & 14 & 0.950 & 0.960 & 0.954 \\
Music      & 20 & 0.965 & 0.951 & \textbf{0.958} & 20 & 0.960 & 0.955 & 0.957 \\
Politics   & 22 & 0.968 & 0.962 & \textbf{0.964} & 22 & 0.945 & 0.961 & 0.953 \\
Science    & 20 & 0.949 & 0.943 & \textbf{0.945} & 20 & 0.922 & 0.933 & 0.926 \\
\midrule
\textbf{Average} & 92 & \textbf{0.970} & \textbf{0.963} & \textbf{0.966}
                 & 92 & 0.938 & 0.951 & 0.944 \\
\bottomrule
\end{tabular}
\end{table}

\subsection{Qualitative Examples}

Table~\ref{tab:examples} shows representative cluster naming examples
from the five domains, illustrating where each method succeeds or struggles.

\begin{table}[h]
\centering
\caption{Selected cluster naming examples. \textcolor{green!60!black}{Green}
= correct or near-correct; \textcolor{red!70!black}{red} = incorrect.}
\label{tab:examples}
\setlength{\tabcolsep}{4pt}
\small
\begin{tabularx}{\textwidth}{l l X X X}
\toprule
\textbf{Domain} & \textbf{Cluster} & \textbf{Ground Truth}
  & \textbf{BERT Name} & \textbf{LLM Name} \\
\midrule
AI         & 11 & conference     & \textcolor{green!60!black}{conference}     & \textcolor{green!60!black}{Conference} \\
AI         & 9  & university     & \textcolor{green!60!black}{university}     & \textcolor{green!60!black}{Universities} \\
AI         & 1  & researcher     & \textcolor{green!60!black}{physicist}      & \textcolor{green!60!black}{Computer Scientist} \\
AI         & 13 & person         & \textcolor{red!70!black}{vegetarian}       & \textcolor{green!60!black}{Person} \\
\midrule
Literature & 2  & writer         & \textcolor{green!60!black}{poet}           & \textcolor{green!60!black}{Author} \\
Literature & 13 & literarygenre  & \textcolor{green!60!black}{genre}          & \textcolor{orange!80!black}{Literary work} \\
Literature & 0  & book           & \textcolor{red!70!black}{vegetarian}       & \textcolor{green!60!black}{Novel} \\
\midrule
Music      & 3  & band           & \textcolor{green!60!black}{band}           & \textcolor{green!60!black}{Band} \\
Music      & 7  & musicgenre     & \textcolor{green!60!black}{genre}          & \textcolor{green!60!black}{Music genres} \\
Music      & 4  & musicalartist  & \textcolor{red!70!black}{vegetarian}       & \textcolor{green!60!black}{Singer} \\
\midrule
Politics   & 7  & politicalparty & \textcolor{green!60!black}{party}          & \textcolor{green!60!black}{Political party} \\
Politics   & 11 & politician     & \textcolor{green!60!black}{democrat}       & \textcolor{green!60!black}{Politician} \\
Politics   & 16 & person         & \textcolor{red!70!black}{vegetarian}       & \textcolor{red!70!black}{Actor} \\
\midrule
Science    & 15 & protein        & \textcolor{green!60!black}{protein}        & \textcolor{green!60!black}{Protein} \\
Science    & 9  & scientist      & \textcolor{green!60!black}{physicist}      & \textcolor{green!60!black}{Physicist} \\
Science    & 5  & astronomicalobj & \textcolor{green!60!black}{planet}        & \textcolor{green!60!black}{Planet} \\
Science    & 14 & person         & \textcolor{red!70!black}{vegetarian}       & \textcolor{red!70!black}{Actor} \\
\bottomrule
\end{tabularx}
\end{table}

% ─── 4. Analysis ─────────────────────────────────────────────────────────────
\section{Analysis}

\subsection{BERT Naming Strengths and Limitations}

BERT naming achieves the highest average F$_1$ (0.966) at zero additional
inference cost since the encoder is already loaded for entity typing.
It excels at domain-specific technical terms (\textit{conference},
\textit{university}, \textit{protein}, \textit{genre}) where BERT's
pretraining vocabulary is well-aligned with the entity types.

However, BERT naming exhibits a known artefact: \textbf{\texttt{vegetarian}}
appears repeatedly across domains as a predicted token. This occurs because
the prompt template \texttt{``X is a [MASK].''} in unusual entity contexts
triggers BERT to predict words from common fill-in-the-blank patterns
unrelated to entity types. This affects clusters 13 (AI), 0 (Literature),
1 and 4 and 19 (Music), 13 and 16 (Politics), and 14 (Science).

An additional limitation is that BERT predicts a \emph{single sub-word
token}, so multi-word entity types such as
\texttt{astronomicalobject}, \texttt{chemicalelement}, or
\texttt{musicalartist} can never be predicted exactly.

\subsection{LLM Naming Strengths and Limitations}

LLM naming (Llama-3.1-8B, $n=16$) produces more natural and
multi-word type names (\textit{Computer Scientist},
\textit{Political party}, \textit{Scientific Journals}), avoiding the
single-token constraint. It never produces nonsensical artefacts like
\texttt{vegetarian}.

However, the LLM occasionally over-generalises:
\begin{itemize}[leftmargin=1.5em]
  \item \textit{Stadiums} for \texttt{country} (AI, cluster 4)
  \item \textit{Actor}/\textit{Celebrity} for \texttt{person} (Politics,
    Science)
  \item \textit{Song} for \texttt{album} (Music, cluster 6)
\end{itemize}
This occurs when the sampled entities are dominated by false positives or
semantically diverse spans that mislead the model.

\subsection{Domain-level Observations}

\begin{itemize}[leftmargin=1.5em]
  \item \textbf{AI} shows the largest gap between methods
    ($\Delta$F$_1 = 0.060$), likely because BERT's pretraining includes
    substantial AI literature where \textit{conference} and
    \textit{university} appear in exactly the fill-in-the-blank pattern.
  \item \textbf{Music} shows the smallest gap ($\Delta$F$_1 = 0.001$),
    with both methods performing well on concrete musical concepts.
  \item \textbf{Science} yields the lowest scores for both methods,
    reflecting the domain's fine-grained types
    (\textit{chemicalelement}, \textit{astronomicalobject},
    \textit{academicjournal}) that require specialised vocabulary.
\end{itemize}

% ─── 5. Conclusion ───────────────────────────────────────────────────────────
\section{Conclusion}

Both naming strategies produce semantically meaningful cluster labels
without any labelled data. BERT naming is the stronger method
on average (F$_1$~0.966 vs 0.944) and comes at no extra cost, while LLM
naming provides more readable, multi-word labels and avoids degenerate
predictions. The two approaches are complementary: BERT naming is
recommended for technical domains with single-word type names, while
LLM naming is preferred when human-readable output or multi-word labels
are required.

Future work could combine both approaches: use BERT naming as a candidate
and LLM naming as a verifier, or fine-tune the LLM on domain-specific
entity samples to reduce hallucination.

% ─── References ──────────────────────────────────────────────────────────────
\begin{thebibliography}{9}

\bibitem{genest2025owner}
P.-Y.~Genest, P.-E.~Portier, E.~Egyed-Zsigmond, and M.~Lovisetto,
``OWNER -- Toward Unsupervised Open-World Named Entity Recognition,''
\textit{IEEE Access}, vol.~13, pp.~50077--50096, 2025.

\bibitem{zhang2020bertscore}
T.~Zhang, V.~Kishore, F.~Wu, K.~Q.~Weinberger, and Y.~Artzi,
``BERTScore: Evaluating Text Generation with BERT,''
in \textit{Proc. ICLR}, 2020.

\end{thebibliography}

\end{document}
